\chapter*{Conclusioni}
\addcontentsline{toc}{chapter}{Conclusioni}

Il presente lavoro ha affrontato il problema della formazione sulla sicurezza hardware in contesti dove l'accesso a laboratori fisici è limitato, proponendo una piattaforma web che simula l'intero flusso di lavoro dell'analisi di un dispositivo embedded --- dall'ispezione visiva del circuito stampato fino all'estrazione e all'analisi del firmware --- in un ambiente completamente virtualizzato e accessibile da browser.

La piattaforma PCB-CTF è stata progettata attorno a due principi guida: la \textbf{fedeltà dell'interazione}, per garantire che l'esperienza dello studente sia trasferibile al contesto reale, e la \textbf{configurabilità dichiarativa}, per consentire ai docenti di creare esercizi personalizzati senza modificare il codice sorgente.

Sul piano dell'interazione, gli strumenti simulati --- multimetro con rumore digitale, adattatore UART con validazione incrociata, programmatore SPI con sei sonde role-based --- replicano la logica operativa delle controparti reali. Il terminale, pur operando su dati predefiniti, riproduce in modo convincente il comportamento di un sistema embedded, dalla sequenza di boot del bootloader U-Boot fino all'esplorazione del filesystem Linux. Il meccanismo di coordinate normalizzate e snap automatico delle sonde rende l'interazione fluida e intuitiva, senza sacrificare il realismo.

Sul piano dell'authoring, il pannello di configurazione consente di definire esercizi arbitrariamente complessi attraverso un editor visuale. Il sistema di sei tipi di output del terminale, combinato con i vincoli contestuali sui comandi e il meccanismo di flag progressive, offre un livello di espressività sufficiente a replicare scenari di analisi realistici, come dimostrato dal caso di studio basato sul router TP-Link WR841N.

Le scelte architetturali si sono rivelate adeguate alla complessità del sistema. La gestione dello stato tramite quattro store Zustand indipendenti ha garantito una separazione chiara delle responsabilità, senza i costi di complessità di un sistema centralizzato. L'assenza di un database e l'adozione di file JSON per la persistenza hanno mantenuto il progetto leggero e immediatamente deployabile, coerentemente con la natura didattica dell'applicazione.

Il caso di studio ha dimostrato che la piattaforma è in grado di supportare un percorso formativo completo, dalla ricognizione hardware all'analisi delle vulnerabilità software, in un formato guidato e gamificato. Lo studente attraversa le stesse fasi concettuali di un'analisi reale --- identificare, misurare, connettersi, estrarre, esplorare, scoprire --- acquisendo competenze metodologiche che costituiscono una base solida per il successivo approccio all'hardware fisico.

Le limitazioni individuate --- l'assenza di esecuzione reale nel terminale, il mancato supporto di protocolli come JTAG e I\textsuperscript{2}C, la mancanza di multi-utenza e persistenza del progresso --- indicano chiaramente le direzioni di evoluzione futura. In particolare, l'integrazione con laboratori remoti per un approccio ibrido simulazione-hardware, e l'introduzione di un sistema di scoring per contesti competitivi, rappresentano le estensioni più promettenti.

In definitiva, PCB-CTF colma un vuoto nell'offerta formativa sulla sicurezza hardware, offrendo un'alternativa accessibile e scalabile ai laboratori fisici. Non si propone di sostituire l'esperienza su hardware reale, ma di preparare lo studente ad affrontarla con le competenze concettuali e la metodologia già acquisite, riducendo la barriera d'ingresso a un campo che, per la crescente pervasività dei dispositivi IoT, è destinato a diventare sempre più rilevante.
